\section{Limitations：从 Benchmark Capability 到 Software Reliability}

课堂最后没有停在 main figure，而是展示 brittle completion、model-size/compute 关系和结论。高 pass@k 能证明在窄 benchmark 上存在功能正确样本，却不能自动覆盖 specification、security、composition 与维护。

\subsection{Brittleness 与 composition}

本节先看 failures，再解释 main figure 的边界。Prompt brittleness 与 composition 退化说明短函数 benchmark 无法代表 repository-scale software。

\lecturefigure{slide-44-limitations.jpg}{Codex limitations 包括 prompt brittleness 与更复杂 composition 上的快速退化。}{官方视频 46:53。}

\begin{knowledgebox}{读图：左边是局部脆弱，右边是规模趋势}
一个小 prompt 变化可能触发错误 completion；随着任务需要组合更多步骤，准确率下降。模型规模能改善曲线，却不保证外推到长程序和真实 repository。应测试 paraphrase、hidden requirements、multi-file dependency、security 和 regression。
\end{knowledgebox}

\begin{warningbox}{Functional correctness 仍不是 production readiness}
生产代码还需要 readability、types、security、performance、licenses、dependency policy、observability 和 human review。即使 unit tests 全过，也可能存在未覆盖输入或危险副作用。模型输出应进入软件工程流程，而不是绕过流程。
\end{warningbox}

\subsection{结论 slide 的四个持久判断}

最后把结论 slide 改写为可检验命题：生成建模、scale、跨模态迁移与 sampling 各自在哪些 data、metric 和 cost 条件下成立。回看整场时应把 sample quality、representation transfer、benchmark pass 与 deployment reliability 放在不同列，避免用一个成功故事覆盖所有层级。

\lecturefigure{slide-45-conclusion.jpg}{讲者用 generative pretraining、scale、cross-modality 与 sampling 总结整场。}{官方视频 47:38。}

\begin{knowledgebox}{读图：把结论改写成可检验命题}
生成建模可利用无标签数据；scale 常提高 sample 与 transfer；同一 recipe 可迁移到 image/code；sampling 能放大分布中的 rare success。每条都需要限定 data、model、metric 与 cost。最重要的延伸是：当 verifier 存在时，生成模型可与搜索形成系统；当 verifier 不可靠时，多采样也可能放大风险。
\end{knowledgebox}

\teachervoice{讲者在结尾同时强调 generative modeling 的广泛性与 limitations。课堂提醒：Codex 的亮点不是“代码问题已解决”，而是 code 提供了一个能执行、能采样、能验证的研究环境，让模型分布和系统搜索可以被分别测量。}

\subsection{本章小结}

Benchmark capability 是必要证据，不是软件可靠性。Codex 让我们看见生成、采样、测试和选择的系统结构，也让 prompt brittleness、composition 与安全边界更加清楚。

